\documentclass[../preamble]{subfiles}
\begin{document}

\section{The Opposite Category}
\label{sec:opposite}
\concept{opposite-category}{Opposite Category}
\depends{category}
\implements{haskell}{Cat.Opposite}
\implements{agda}{Cat.Opposite}
\implements{lisp}{opposite}

\begin{definition}[Opposite category]
  \label{def:opposite-category}
  Let $\Category{C}$ be a category.
  The \newterm{opposite category} (or \newterm{dual category})
  $\newmath{\Category{C}^{\mathrm{op}}}$ is defined by the following data
  (\stacksref{001M}; \nlabref{opposite+category}):
  \begin{enumerate}
    \item $\Ob(\Category{C}^{\mathrm{op}}) = \Ob(\Category{C})$.
    \item For each pair of objects $X, Y$,
      \begin{equation}
        \label{eq:op-hom}
        \Hom_{\Category{C}^{\mathrm{op}}}(X, Y)
        = \Hom_{\Category{C}}(Y, X).
      \end{equation}
    \item For each object $X$,
      the identity morphism is $\newmath{\id^{\mathrm{op}}_X} = \id_X$.
    \item For morphisms
      $f \in \Hom_{\Category{C}^{\mathrm{op}}}(X, Y)$
      and $g \in \Hom_{\Category{C}^{\mathrm{op}}}(Y, Z)$
      (i.e.\ $f \colon Y \to X$ and $g \colon Z \to Y$ in $\Category{C}$),
      the \newterm{opposite composition} is
      \begin{equation}
        \label{eq:op-comp}
        g \circ^{\mathrm{op}} f = f \circ g \in \Hom_{\Category{C}}(Z, X)
        = \Hom_{\Category{C}^{\mathrm{op}}}(X, Z).
      \end{equation}
  \end{enumerate}
\end{definition}

\begin{remark}[Opposite category satisfies the axioms]
  \label{rem:opposite-axioms}
  We verify that $\Category{C}^{\mathrm{op}}$ is indeed a category.
  \begin{itemize}
    \item \textbf{Associativity.}
      For morphisms $f \colon X \to Y$, $g \colon Y \to Z$, $h \colon Z \to W$
      in $\Category{C}^{\mathrm{op}}$
      (i.e.\ $f \colon Y \to X$, $g \colon Z \to Y$, $h \colon W \to Z$ in $\Category{C}$),
      \[
        h \circ^{\mathrm{op}} (g \circ^{\mathrm{op}} f)
        = h \circ^{\mathrm{op}} (f \circ g)
        = (f \circ g) \circ h
        = f \circ (g \circ h)
        = (g \circ h) \circ^{\mathrm{op}} f
        = (h \circ^{\mathrm{op}} g) \circ^{\mathrm{op}} f,
      \]
      using associativity of composition in $\Category{C}$.

    \item \textbf{Identity.}
      For $f \colon X \to Y$ in $\Category{C}^{\mathrm{op}}$
      (i.e.\ $f \colon Y \to X$ in $\Category{C}$),
      \[
        f \circ^{\mathrm{op}} \id^{\mathrm{op}}_X
        = \id_X \circ f = f
        \quad \text{and} \quad
        \id^{\mathrm{op}}_Y \circ^{\mathrm{op}} f
        = f \circ \id_Y = f,
      \]
      using the identity laws in $\Category{C}$.
  \end{itemize}
\end{remark}

\begin{proposition}
  \label{prop:opposite-involution}
  For any category $\Category{C}$, there is an equality
  $(\Category{C}^{\mathrm{op}})^{\mathrm{op}} = \Category{C}$.
\end{proposition}

\begin{proof}
  The objects and hom-sets of $(\Category{C}^{\mathrm{op}})^{\mathrm{op}}$
  are recovered by the double reversal
  \[
    \Hom_{(\Category{C}^{\mathrm{op}})^{\mathrm{op}}}(X, Y)
    = \Hom_{\Category{C}^{\mathrm{op}}}(Y, X)
    = \Hom_{\Category{C}}(X, Y),
  \]
  and composition $f \circ^{\mathrm{op\,op}} g = g \circ^{\mathrm{op}} f = f \circ g$
  agrees with composition in $\Category{C}$.
\end{proof}

\concept{contravariant-functor}{Contravariant Functor}
\depends{opposite-category}
\depends{functor}

\begin{definition}[Contravariant functor]
  \label{def:contravariant-functor}
  A \newterm{contravariant functor} from $\Category{C}$ to $\Category{D}$
  is a \newterm{(covariant) functor}
  $\Functor{F} \colon \Category{C}^{\mathrm{op}} \lto \Category{D}$.

  Explicitly, $\Functor{F}$ assigns to each object $X \in \Ob(\Category{C})$
  an object $\Functor{F}(X) \in \Ob(\Category{D})$,
  and to each morphism $f \colon X \to Y$ in $\Category{C}$
  a morphism
  \begin{equation}
    \label{eq:contravariant-hom}
    \begin{aligned}
      \Functor{F} \colon \Hom_{\Category{C}}(X, Y)
        &\lto \Hom_{\Category{D}}(\Functor{F}(Y), \Functor{F}(X)) \\
      f &\lmapsto \Functor{F}(f)
    \end{aligned}
  \end{equation}
  in $\Category{D}$,
  subject to $\Functor{F}(\id_X) = \id_{\Functor{F}(X)}$
  and $\Functor{F}(g \circ f) = \Functor{F}(f) \circ \Functor{F}(g)$.
\end{definition}

\begin{remark}
  The representable presheaf functor
  \begin{equation}
    \begin{aligned}
      \Hom_{\Category{C}}(-, Y) \colon \Category{C}^{\mathrm{op}}
        &\lto \Set \\
      X &\lmapsto \Hom_{\Category{C}}(X, Y)
    \end{aligned}
  \end{equation}
  is the canonical example of a contravariant functor.
  The collection of all such functors organises into the presheaf category
  $\newmath{[\Category{C}^{\mathrm{op}}, \Set]}$,
  into which $\Category{C}$ embeds fully faithfully
  via the Yoneda embedding.
\end{remark}

\end{document}
